% Figure: transfer characteristic of a class-B push-pull output stage, showing
% the dead band near the origin (crossover distortion) where neither transistor
% conducts, and the class-AB curve that a small bias current straightens. Plain
% pgfplots, output-versus-input voltage. No circuitikz.
\begin{figure}[htbp]
\centering
\begin{tikzpicture}
\begin{axis}[
  width=10.5cm, height=6.0cm,
  xlabel={input voltage $v_{\text{in}}$ (V)},
  ylabel={output voltage $v_{\text{out}}$ (V)},
  xmin=-3, xmax=3, ymin=-3, ymax=3,
  axis lines=middle,
  tick label style={font=\tiny},
  label style={font=\scriptsize},
  xtick={-2,-1,1,2}, ytick={-2,-1,1,2},
  clip=false,
]
% class B: dead band |v_in| < 0.6 V (one V_BE), unity slope outside
\addplot[engred, very thick, samples=200, domain=-3:3]
  {(x > 0.6)*(x-0.6) + (x < -0.6)*(x+0.6)};
% class AB / ideal: straight unity line for comparison
\addplot[engblue, thick, dashed, domain=-3:3, samples=2] {x};
\node[engred, font=\tiny, anchor=west] at (axis cs:1.1,0.25) {class B};
\node[engblue, font=\tiny, anchor=south east] at (axis cs:2.6,2.6) {class AB};
% mark the dead band
\draw[enggray, dotted] (axis cs:-0.6,0) -- (axis cs:-0.6,-0.6);
\draw[enggray, dotted] (axis cs:0.6,0) -- (axis cs:0.6,0.6);
\node[enggray, font=\tiny, anchor=north] at (axis cs:0,-0.15)
  {dead band $\pm V_{BE}$};
\end{axis}
\end{tikzpicture}
\caption[Crossover distortion of a class-B output stage]{The transfer
characteristic of a class-B push-pull output stage (red) is flat near the
origin: within one base-emitter drop of zero input, neither the upper nor the
lower transistor conducts, so the output sits at zero and a small signal
disappears into the dead band, the crossover distortion that a class-B stage
plants on every zero crossing. A small quiescent bias current keeps both
devices barely on through the crossover, straightening the curve toward the
ideal unity line (blue dashed) and turning the stage into class AB at the cost
of a standing dissipation.}
\label{fig:vol04ch09:class-b-crossover}
\end{figure}
